\clearpage
\pdfbookmark[0]{Abstract}{abstract}
\chapter*{ABSTRACT}
\addcontentsline{toc}{chapter}{ABSTRACT}

This study compares cooperation structures among information channels that emerge after classifiers are trained using samples from two different sources. In the reference condition, classifiers are trained on real samples; in the synthetic conditions, the same classifiers are trained on samples produced by probabilistic models parameterized from the real training data. In each condition, the cooperation structure is characterized by how the ranking of the best channel subsets changes as the training-set size increases. The study then investigates to what extent the structure obtained under synthetic training reproduces the structure obtained under real training. The protocol compares three training conditions---real data, a class-conditional single Gaussian, and a Gaussian mixture model---evaluated on the same fixed real test set. Every non-empty channel subset is examined over a grid of per-class sample sizes through Monte Carlo simulation with incrementally nested samples. Experiments use the Occupancy Detection, UCI Air Quality, and SUPPORT2 datasets with linear SVM, Logistic Regression, Gaussian Naive Bayes, and, additionally for SUPPORT2, Random Forest. Preservation is evaluated through ranking correlation, pairwise \WinnersAgreement{}, and progressive similarity of $N^{\star}$ matrices, defined by the last observed crossing between loss curves. Results show that structural fidelity is multidimensional: rankings and pairwise winner relations can be strongly preserved even when $N^{\star}$ similarity is moderate or low. In Air Quality, Gaussian Naive Bayes nearly perfectly preserved rankings and winners. In SUPPORT2, Random Forest combined poor absolute predictive performance with the highest $N^{\star}$ similarity among the classifiers evaluated in that scenario, despite thousands of defined crossings. Training on samples from the single Gaussian was competitive with or superior to training on GMM samples in some metrics, although the GMM visually reproduced some geometries more closely. Predictive performance, geometric fidelity, and preservation of cooperation are therefore distinct properties. The study does not demonstrate cost reduction, but it provides structural objects for future cost-efficiency analysis and for investigating projections to larger sample-size regimes.

\vspace{\onelineskip}
\noindent\textbf{Keywords}: information channels; cooperative classification; synthetic data; subset selection; Monte Carlo simulation; structural fidelity.
